\section{Дискретизация фазы}
\label{sec:descent-phase-disc}
\label{sec:coords-internal}

На макроэтаже волновая функция допускает гидродинамическую запись Маделунга,
\[
  \wavePsi = \sqrt{\massDensity}\,e^{i\phaseVarphi/\planckHbar},
\]
где $\massDensity = |\wavePsi|^2$ есть плотность, а безразмерный показатель экспоненты $\varphi = \phaseVarphi/\planckHbar$ измеряет действие в единицах~$\planckHbar$.
На одном узле~$\caPlanckCell$ ту же роль играют сумма квадратов модулей $|z|^2$ и аргументы компонент спинора~$z_k$.

Однако на одной планковской ячейке фаза не может принимать произвольное вещественное значение из непрерывной окружности.
Если бы аргумент компоненты~$z_k$ задавался бесконечной десятичной дробью, один-единственный узел хранил бы бесконечное количество информации, что прямо противоречит уже установленному пределу Бекенштейна \eqref{eq:bekenstein-bits}.
Следовательно, окружность фаз на узле обязана замкнуться в \emph{конечное циклическое кольцо}.

Определим одновременно и полное число делений этого кольца, и величину минимального шага, который на нём откладывается.
В квантовой механике (\cref{sec:qm}) мы уже видели, что проекция спина электрона \eqref{eq:spin-z} и соотношение неопределённостей Гейзенберга $\Delta S \ge \planckHbar/2$ выделяют неразложимый квант действия, равный половине постоянной Планка:
\[
  \ladderAction = \frac{\planckHbar}{2}.
\]
Деля этот минимальный квант действия на~$\planckHbar$, получаем наименьший физически различимый сдвиг фазы за один элементарный акт переноса:
\[
  \symCapitalDelta\symPhi_{\min} = \frac{\ladderAction}{\planckHbar} = \frac{1}{2}\ \text{рад}.
\]

\begin{theorem}[Информационный бюджет элементарной ячейки]
\label{thm:cell-bit-budget}
Пусть элементарная ячейка~$\hv$ имеет линейный размер $R\sim\caStepSpace$ и объём $v_{\hv}\sim\caStepSpace^3$.
Тогда предел Бекенштейна (\cref{def:bekenstein-bound}) и условие гравитационного не-схлопывания \eqref{eq:energy-bound} фиксируют предельную информационную ёмкость одной ячейки:
\[
  B_{\hv} = \frac{2\symPi}{\ln 2}\approx 9{,}0647\ \text{бита},\qquad
  \phaseRingCount = 2^{\lfloor B_{\hv}\rfloor} = 2^{9} = 512,
\]
\[
  B_{\hv} = \log_2 \phaseRingCount + \varepsilon_B,\qquad \varepsilon_B \approx 0{,}065.
\]
На этом кольце $\mathbb{Z}_{\phaseRingCount}$ из $512$ делений минимальный фазовый шаг Гейзенберга $\symCapitalDelta\symPhi_{\min}=\tfrac{1}{2}$ рад укладывается в полном обороте~$2\symPi$ в виде
\[
  \phaseSectorCount = \left\lceil \frac{2\symPi}{\symCapitalDelta\symPhi_{\min}}\right\rceil = \lceil 4\symPi\rceil = 13\ \text{секторов},\qquad
  \symCapitalDelta\symPhi_{\mathrm{disc}}
  = \left\lceil \phaseRingCount\,\frac{\symCapitalDelta\symPhi_{\min}}{2\symPi}\right\rceil = 41\ \text{тик кольца}.
\]
\end{theorem}

\begin{proof}
\pfrom{предела Бекенштейна (\cref{def:bekenstein-bound}) при радиусе области $R=\caStepSpace$ и энергии $E$}
число битов в ячейке ограничено величиной
\[
  B_{\hv} \le \frac{2\symPi R E}{\planckHbar\,\speedC\,\ln 2}
\]
(\eqref{eq:bekenstein-bits}).
\pusing{потолок энергии на масштабе ячейки $E = E_P = \planckHbar\,\speedC/\caStepSpace$ (\eqref{eq:planck-energy})}
находим точное насыщение границы на одной планковской ячейке:
\psubst{$R=\caStepSpace$ и $E=E_P$}
\[
  B_{\hv} = \frac{2\symPi}{\ln 2}.
\]
\pfollows
целое число полных двоичных разрядов регистра ячейки равно
\[
  \lfloor B_{\hv}\rfloor = 9,\qquad
  \phaseRingCount = 2^{9} = 512,\qquad
  \varepsilon_B = B_{\hv} - \lfloor B_{\hv}\rfloor \approx 0{,}065.
\]

\smallskip
\emph{Разметка фазового кольца.}
\pfrom{минимального фазового шага $\symCapitalDelta\symPhi_{\min}=\tfrac{1}{2}$ рад}
полный оборот $2\symPi$ вмещает $\lceil 2\symPi/\symCapitalDelta\symPhi_{\min}\rceil = \lceil 4\symPi\rceil = 13$ секторов.
\psubst{$\phaseRingCount = 512$}
на один минимальный шаг Гейзенберга приходится целое число делений кольца:
\[
  \symCapitalDelta\symPhi_{\mathrm{disc}}
  = \left\lceil \phaseRingCount\,\frac{\symCapitalDelta\symPhi_{\min}}{2\symPi}\right\rceil = \left\lceil \frac{512}{4\symPi}\right\rceil = 41.
\]
\end{proof}

Читатель, привыкший к гладким производным, в этом месте обычно чувствует, что лесенка под ногой внезапно кончилась.
Если вы слегка ушиблись при переходе от непрерывной фазы~$\phaseVarphi\in\spaceR$ к целым вычетам по модулю~$512$ и шагу~$41$ --- простите: дальше ступеньки снова пойдут ровно, но сам этот обрыв здесь неизбежен, потому что непрерывная прямая на масштабе одной ячейки действительно кончается.

За один такт~$\htick$ фазовый счётчик пробегает кольцо $\spaceZ_{\phaseRingCount}$ с элементарным поворотом $\symOmegaLower = e^{i2\symPi/\phaseRingCount}$, так что угловое положение на окружности~$S^1$ имеет вид $\symOmegaLower^{\phaseVarphi_{\mathrm{disc}}}$.

\begin{remark}[Различие между координатой фазы $\phaseVarphi_{\mathrm{disc}}$, генератором $\symOmegaLower$ и импульсом $\symCapitalPhi$]
\label{rem:phi-omega-Phi}
Индекс $\phaseVarphi_{\mathrm{disc}}\in\spaceZ_{\phaseRingCount}$ есть целочисленная \emph{координата} фазы на кольце.
Множитель $\symOmegaLower$ --- фиксированный корень степени~$512$ из единицы, задающий перевод индекса в поворот плоскости.
Величина $\symCapitalPhi\in\spaceZ_{\phaseRingCount}$ обозначает \emph{приращение фазы (импульс) за один такт}, поворачивающее спинор на $R(\symCapitalPhi)=\symOmegaLower^{\symCapitalPhi}$.
При переходе к макроописанию~$\layerMacro$ дискретный индекс переходит в угол $\phaseVarphi = \phaseVarphi_{\mathrm{disc}}\cdot(2\symPi/\phaseRingCount)$.
\end{remark}

Из девятки полных разрядов ячейки ($\lfloor B_{\hv}\rfloor = 9$) три разряда обслуживают трёхмерное пространственное ветвление и знак, а оставшиеся $\mathrm{frac\_bits} = 9 - 3 = 6$ разрядов задают двоичную сетку амплитуды с шагом $z_{\min} = 2^{-6} = 1/64$ ($\symMu\in\spaceZ_{64}$).
В результате каждая компонента спинора на узле записывается как элемент конечного гауссова кольца
\[
  Z = U + iV \in \spaceZ_{\phaseRingCount}[i].
\]
Если теперь разделить текущую фазу $\phaseVarphi_{\mathrm{disc}}\in\{0,\ldots,511\}$ с остатком на порог Гейзенберга $\symCapitalDelta\symPhi_{\mathrm{disc}} = 41$:
\[
  \phaseVarphi_{\mathrm{disc}} \equiv k_{\symPhi}\cdot\symCapitalDelta\symPhi_{\mathrm{disc}} + \phaseVarphi_{\mathrm{f}}
  \pmod{\phaseRingCount},
\]
то частное $k_{\symPhi}\in\spaceZ_{13}$ покажет, сколько полных квантов действия~$\ladderAction$ уже набрано в узле, а остаток $\phaseVarphi_{\mathrm{f}}\in\{0,\ldots,40\}$ задаст подпороговую фазу внутри одного сектора.

\begin{proposition}[Каноническое разложение фазы и шов замыкания окружности]
\label{prop:canonical-phase}
Всякому состоянию $\phaseVarphi_{\mathrm{disc}}\in\spaceZ_{\phaseRingCount}$ однозначно отвечает каноническая пара $(k_{\symPhi},\phaseVarphi_{\mathrm{f}})$ с наименьшим неотрицательным номером сектора~$k_{\symPhi}$.
Поскольку $13\times 41 = 533 > 512$, тринадцатый сектор при замыкании полного круга перекрывается с началом отсчёта ровно на $21$ деление:
\[
  \phaseRingCount = \phaseSectorCount\cdot\symCapitalDelta\symPhi_{\mathrm{disc}} - (\phaseSectorCount + \vevHierarchySteps) = 13\cdot 41 - 21 = 512.
\]
\end{proposition}

Этот целочисленный избыток $533 - 512 = 21 = \phaseSectorCount + \vevHierarchySteps$ (сумма размера одетого кластера~$13$ и глубины битовой иерархии~$8$) играет фундаментальную роль на верхних этажах: он равен числу спиновых пар на второй координационной сфере кубооктаэдра.
Отношение же одного шага Гейзенберга к полному кольцу $\symCapitalDelta\symPhi_{\mathrm{disc}}/\phaseRingCount = 41/512$ представляет собой дискретную запись фазового вычета одного такта $r = 1/(4\symPi)$ и, как показано в главе~\ref{ch:alpha}, отлично от постоянной тонкой структуры~$\alphaFs$.
